% !TEX root = chapter-standalone.tex



\section{Uses of a functor: more examples}
\label{sec:more-examples-functors}


\devel{

\subsection{Here are notes from the discussion on 29.06.2026}


The possible jobs of a functor:
\begin{enumerate}
\item embedding-like (also deifnitng structrure of A from a map A → B where B has structure). the functor defines the structure. you choose nay function, you assert that it must be fucntor

\item compiling-like, projection like. You need to find a functor (loses infomration, but preserves the “Structure”)

example: a compiler source code → machine code. 

example: the deployment computation grah → resource graph

\item there is a function, and a property of it is encoded by the enforcing of it being a functor

example:  monotone functions 

\item the functor describes a lifting/generaliation -  a recipe / construatction

\item the functors encodes constraints that something must respect. (the functor is there before the categroy). spivak: the functor encodes integrity constrinats on the data - you REJECT data that is not compatible; this is generalize in a sense to enocoding diagrams via functors 

\item the search for a functor: the solution of a problem consists of specifying a functor
(example: berg)

\end{enumerate}




\subsection{Functors---operational clusters}


(Small excerpt from a Claude session collecting many examples and identifying use patterns)


\begin{enumerate}
\item[F1.] \textbf{Instantiation: the functor is the model.}
Fix a small category as a template (schema, shape, time axis, alphabet); a functor out of it into $\Set$ fills the template with content.
One studies the functor itself as a single object: a system, an instance.
Members: databases, $G$-actions, Moore actions, graphs-as-functors, dynamical systems from a time monoid, family trees, Markov evolution, edge labellings, the amounts functor, the behaviors presheaf.
Signature: small/syntactic domain, $\Set$-like codomain.
The domain's equations become the model's integrity constraints.

\item[F2.] \textbf{Probing: picking out patterns.}
A functor from a tiny ``walking'' category into an arbitrary one selects a configuration.
Members: diagrams-as-functors, the walking isomorphism/idempotent, Perk, planning-as-search.

\item[F3.] \textbf{Batch construction: containers and free structures.}
Apply a construction to every object at once, with transport along morphisms (``map'').
Containers: powerset, $\mathcal{D}$, list, Maybe, multiset, polynomial functors, $A \times (-)$, $(-)^A$.
Free: free vector space, free category, discrete/codiscrete.
Signature: endofunctor-like; functoriality is the API guarantee.

\item[F4.] \textbf{Extraction: invariants, behaviors, forgetting.}
Map a rich category to a poorer one, keeping only what matters---to summarize, or to obstruct (unequal images prove non-isomorphism).
Members: $\pi_0$, fixed points, Perk, cycle-space, the rate functor, programs-to-functions, behavior functors, quotient-set, metric-to-preorder, clustering, forgetful functors.
Many are ``quotient, functorially''.

\item[F5.] \textbf{Translation: dictionaries between formalisms.}
Move a whole problem to a comparable formalism where it's easier.
Members: log of rates, partial functions vs.\ pointed sets, $\FinSet\op$ vs.\ finite Boolean algebras, the beam-model formalization, machine-to-circuit compilation.
Faithfulness/fullness make the translation trustworthy.

\item[F6.] \textbf{Pullback: contravariant observation.}
Data flows backwards along system maps by precomposition.
Members: observables $V^X$, preimage powerset, the behaviors presheaf, data migration, Legendre--Fenchel.
The operational reason $\CatCop$ exists.

\item[F7.] \textbf{Parameter sweep: coherent families over a poset.}
Run a construction at every value of a parameter and record the comparison maps.
Members: clustering across scales, persistence, Markov/Chapman--Kolmogorov, interval-arithmetic isotonicity, the beam-model hierarchy.
Functoriality out of a poset = the sweep is coherent.

\item[F8.] \textbf{Certification: functoriality as a correctness property.}
The operation exists in practice; the act is to verify or refute functoriality.
Home of the non-examples.
Pass: map, single-linkage clustering (average/complete fail), interval extension.
Fail: cheapest-route laxity, entrywise nonlinear matrix ops, $\Aut$, group center.
The cheapest-route case fails equality but satisfies an inequality---a preview of ``lax''.

\item[F9.] \textbf{Plumbing: infrastructure inside \Category.}
Functors as wiring for later machinery: diagonal, projections, slice post-composition, $\Arrow(F)$, inclusions, identities.

\item[F10.] \textbf{The functor as the unknown.}
Problems posed as ``find a functor'' or ``does one exist?'': planning-as-search, database-instance design, Church--Turing as fullness.
\end{enumerate}

\emph{Meta-observation.}
The tags track the shape of domain and codomain, like a type discipline:
small/syntactic domain gives instantiation (F1) or probing (F2); endofunctors give construction (F3); a poor codomain gives extraction (F4) or sweep (F7); an opposite-category domain gives observation (F6); peer-to-peer gives translation (F5).
Announcing a functor's tag may orient students faster than its formal data does.

}


\devel{

\begin{example}[Equivalence Classes]
    \label{exa:equivalence-classes-functor}
    Consider the category $\EquivRel$ from \cref{def:cat-equivalence-relations}.
    We define a functor
    \begin{equation}
        \funa \colon \EquivRel \fto \Set
    \end{equation}
    as follows.

    Given an object $\tup{\setA, \equival_\setA}$ of $\EquivRel$ (that is, a set equipped with an equivalence relation), we define $\funaob(\tup{\setA, \equival_\setA}) = \setA/_{\equival_\setA}$ to be the quotient of $\setA$ by $\equival_\setA$.

    Recall that $\setA/_{\equival_\setA}$ is the set of equivalence classes of $\equival_\setA$.
    Its elements are the subsets of $\setA$ which form the partition of $\setA$ induced by $\equival_\setA$.
    Each such subset is the set of all elements of $\setA$ which are mutually equivalent to each other according to $\equival_\setA$.
    For any element $\ela \setin \setA$, the equivalence class it belongs to is denoted $[\ela]$ and in this case $\ela$ is called a \emph{representative} of the equivalence class $[\ela]$.

    To define $\funamor$, let $\mora \colon \tup{\setA, \equival_\setA} \mto \tup{\setB, \equival_\setB}$ be a morphism in $\EquivRel$.
    We let
    \begin{equation}
        \funamor(\mora) \colon \setA/_{\equival_\setA} \mto_\Set \setB/_{\equival_\setB}, \ [\ela] \mapsto [\mora(\ela)].
    \end{equation}
    It may be readily checked that this function is well-defined, irrespective of the (arbitrary) choice of element $\ela$ used to represent a given equivalence class $[\ela]$.
\end{example}

\todotext{Give an illustration of a possible application of the above, for example using rounding of currency...}

\begin{exercise}
    Prove that the functor defined in \cref{exa:equivalence-classes-functor} is in fact a functor.
\end{exercise}
\begin{solution}
    \todotext{Write this solution}
\end{solution}


\begin{gradedexercise}[\exname{DoubleDualFunctor}]
    \label{ex:DoubleDualFunctor}
    Prove that $\funa$ as defined in \cref{exa:double-dual-functor} is in fact a \SY{functor}.
\end{gradedexercise}
\solutionof{DoubleDualFunctor}

% 2026-11-10: the example on group actions as functors (exa:G-actions-as-functors) and its exercise
% are now covered by \cref{rem:group-actions-as-functors} in the subsection ``Functors generalize monoid actions''.

%\begin{gradedexercise}[\exname{GActionsAsFunctors}]
%    \label{ex:GActionsAsFunctors}
% 
%\end{gradedexercise}
%\solutionof{GActionsAsFunctors}


% 2026-11-10: the graded exercises MultiplicationWithASet, ExponentiationWithASet, CartProdAsFunctor,
% ProbabilityFunctor, DifferentiationFunctor and FixedPointFunctor have been turned into examples (each followed by a short graded
% exercise with the same name and label) at the end of the section ``Functors'' in 10_functors.tex.

\begin{gradedexercise}[\exname{GraphsViaFunctors}]
    \label{ex:GraphsViaFunctors}
    Consider the following category, which has two objects,~$\styleobj{V}$ and~$\styleobj{A}$, and four morphisms: besides the \SY{identity morphisms}, there are two morphisms,~$\stylemorph{s}$ and~$\stylemorph{t}$, from~$\styleobj{A}$ to~$\styleobj{V}$.
    See \cref{fig:graph-cat}.
    Call this category~$\CatSymbol{G}$.

    Can you explain the following statement?
    ``Specifying a \SY{functor}~$\CatSymbol{G} \fto \Set$ is the ``same thing'' as specifying a directed graph''.
\end{gradedexercise}
\solutionof{GraphsViaFunctors}

\begin{marginfigure}
    \centering
    \includesag{graph-cat}
    \caption{}
    \label{fig:graph-cat}
\end{marginfigure}



\begin{example}[Double dual]
    \label{exa:double-dual-functor}
    Let~$\VectRs$ be the category whose objects are all real \SY{vector spaces} and whose morphisms are~\reals-linear maps.
    Composition is the usual composition of linear maps.

    There is an \SY{endofunctor}~$\funa \colon \VectRs \mto \VectRs$ whose action on objects is
    \begin{equation}
        \funa ( \vecspB ) = \vecspB^{**}.
    \end{equation}
    (Recall that~$\vecspB^{**} = \makeset{ \text{linear maps } \vecspB^* \mto \reals } =  \HomSet{\VectRs}{\vecspB^*}{\reals}$).
    The action of~$\funa$ on morphisms is as follows.
    Given a linear map~$\mora \colon \vecspB \mto \vecspC$, we can write
    \begin{equation}
        \defmapperiod{\funa ( \mora ) }{\vecspB^{**}}{\mto}{\vecspC^{**}}{\xi}{[ l \mapsto \xi (f \mthen l)]}
    \end{equation}
\end{example}



    \subsection{The elevation functor}

    \begin{figure}
        \includegraphics[width=8cm]{berg_elevation}
    \end{figure}

}

\subsection{Planning as the search for a functor}
\label{sec:planning-as-functor}

\begin{example}
    \label{exa:planning-as-search-functor}
    Recall the category \Berg introduced in \cref{sec:trekking} and define a category \Plans where objects are activities and morphisms describe activities order constraints, illustrated in \cref{fig:plans_functor} (left).

    For instance, there is a morphism from ``mountain lodge'' to ``panoramic lake'', which describes the plan of going from the lodge area to the lake area.
    We call such morphisms \emph{plans}.
    Plans can be composed via concatenation.
    For instance, given a plan to go from ``mountain lodge'' to ``panoramic lake'', and a plan to go from ``panoramic lake'' to ``peak'', their composition is the plan of going from ``mountain lodge'' to the ``peak'', passing through the ``panoramic lake''.

    When we talk about ``planning'' in this context, we refer to the action of finding a \SY{functor}~$\funa$ from \Plans to \Berg.
    % Let's look at this in more detail.
    The objects of \Berg are tuples~$\tup{p,v}$, where~$p$ represent coordinates of a specific location and~$v\setin \reals^3$ represents velocities.
    Morphisms in \Berg are paths that connect locations.
    For the sake of our planning, we can identify areas of the mountain as sets of locations.
    Such areas are, for instance, the ``mountain lodge'', ``panoramic lake'', and the
    ``peak'' (note that the ``peak'' represents an area corresponding to a single location).
    Given some plans as in \cref{fig:plans_functor} (left), we want to find a map~$P$ which maps each object in \Plans (activity) to an object of \Berg (specific location and velocity).
    Similarly, it must map each morphism in \Plans (activity order constraints) to a morphism in \Berg (specific paths).
    This is illustrated in \cref{fig:plans_functor}.
\end{example}

\begin{figure*}[h!]
    \centering
    \includesag{planning_berg}
    \caption{Planning \SY{functor}.}
    \label{fig:plans_functor}
\end{figure*}

\vfill
\clearpage
\vfill

\todojira{484}{\alphubel: @Gioele: Show that there is a \SY{functor} from reals to cat, which maps an alpha to the category bergalpha}
